\begin{abstract}
Claims about what fine-tuning does to a language model's weights often rest on
the rank of the update $\dW = W_{\text{after}} - W_{\text{before}}$. We measure
four rank measures on real updates and on random matrices of the same shape.
Three measures in common use, the number of singular values above
$0.01\,\sigma_1$, the rank needed for 90\% of the energy, and the entropy
effective rank, put real updates only 0.4--31\% below the noise value, so that
read on their own they describe the updates as close to full rank. Stable rank,
$\|\dW\|_F^2/\|\dW\|_2^2$, puts the same updates 1.6--9$\times$ below noise.
All four detect structure; only stable rank makes its size visible. We show
this on full SFT, DPO and token-matched non-alignment fine-tunes of
Qwen2.5-0.5B and SmolLM2-360M, with seed repeats and layer-bootstrap
intervals, and on seven publicly released base/instruct pairs from 135M to 8B
parameters. The updates combine a few dominant directions with a long tail
that holds most of the energy, and they are at least as concentrated as the
pretrained weights they are added to. The concentration is not specific to alignment:
fine-tuning on Wikipedia text or Python code is as concentrated as instruction
tuning. As a case study we revisit Timber, which cites an unchanged entropy
effective rank between base and instruct weights as evidence that
post-training is superficial. A random perturbation of the same norm leaves
the effective rank just as unchanged. Timber's eRank cutoff for editing $\dW$
leaves 99\% of the update's energy untouched; a stable-rank cutoff that scales
53 times more of it gives the same GSM8K pass@$k$, and removing that tail
outright costs 4 points of pass@1.
\end{abstract}
